%---------------------------------------------------------------------
% UNIT 8 --- PARTIAL DERIVATIVES
% Source: Section 8.9 Self Assessment Questions, pp. 183-184.
%---------------------------------------------------------------------
\chapter{Partial Derivatives}

\srcnote{Source: \emph{Business Mathematics} 5405/1429, \S8.9 \saq{Self Assessment Questions}, pp.~183--184.}

\begin{keyidea}[Hold the other variable constant]
To find $f_x$, differentiate with respect to $x$ and treat $y$ as a
\emph{number}. To find $f_y$, treat $x$ as a number. Every ordinary rule ---
power, product, quotient, chain --- carries over unchanged.

\medskip
For mixed partials, \textbf{Clairaut's theorem} says the order of
differentiation does not matter for the functions in this course:
\[
f_{xy}=f_{yx},\qquad f_{xyz}=f_{yzx}=f_{xzy}=\dots
\]
That is a free correctness check, and Question 3 is built on it.

\medskip
\textbf{Second-derivative test} for $z=f(x,y)$ at a critical point where
$f_x=f_y=0$: with $D=f_{xx}f_{yy}-(f_{xy})^{2}$,
\begin{center}
\small\sffamily
\begin{tabular}{@{}L{4.4cm} L{4.4cm} L{4.4cm}@{}}
\toprule
$D>0$ and $f_{xx}>0$ & $D>0$ and $f_{xx}<0$ & $D<0$ \\
relative minimum & relative maximum & saddle point \\
\bottomrule
\end{tabular}
\end{center}
\noindent
and $D=0$ is inconclusive --- the test simply cannot decide.
\end{keyidea}

%---------------------------------------------------------------------
\section{First partial derivatives}

\begin{question}{Q.1}
Find $f_x$, $f_y$, $f_x(3,4)$ and $f_y(1,2)$ for
\textbf{(a)} $f(x,y)=\dfrac{x+3xy^{2}}{x^{3}+y^{2}}$ \quad
\textbf{(b)} $f(x,y)=y\,e^{5x+2y^{2}}$ \quad
\textbf{(c)} $f(x,y)=(x+1)^{2}+(y-3)^{3}+5xy^{3}-2$
\end{question}

\begin{solution}
\textbf{(a)} Quotient rule, with $y$ held constant for $f_x$.
\[
u=x+3xy^{2}=x(1+3y^{2}),\qquad u_x=1+3y^{2};\qquad
v=x^{3}+y^{2},\qquad v_x=3x^{2}.
\]
\[
f_x=\frac{(1+3y^{2})(x^{3}+y^{2})-x(1+3y^{2})(3x^{2})}{(x^{3}+y^{2})^{2}} .
\]
The factor $(1+3y^{2})$ is common to both terms of the numerator --- taking it
out collapses the whole expression:
\[
f_x=\frac{(1+3y^{2})\bigl[(x^{3}+y^{2})-3x^{3}\bigr]}{(x^{3}+y^{2})^{2}}
=\frac{(1+3y^{2})(y^{2}-2x^{3})}{(x^{3}+y^{2})^{2}} .
\]
\[
f_x(3,4)=\frac{(1+48)(16-54)}{(27+16)^{2}}=\frac{49(-38)}{1849}
=-\frac{1862}{1849}\approx -1.0070 .
\]

Now $f_y$, with $x$ held constant. Here $u_y=6xy$ and $v_y=2y$:
\[
f_y=\frac{6xy(x^{3}+y^{2})-x(1+3y^{2})(2y)}{(x^{3}+y^{2})^{2}}
=\frac{2xy\bigl[3(x^{3}+y^{2})-(1+3y^{2})\bigr]}{(x^{3}+y^{2})^{2}},
\]
and the $3y^{2}$ terms cancel inside the bracket:
\[
f_y=\frac{2xy(3x^{3}-1)}{(x^{3}+y^{2})^{2}} .
\]
\[
f_y(1,2)=\frac{2(1)(2)(3-1)}{(1+4)^{2}}=\frac{8}{25}=0.32 .
\]

\medskip
\textbf{(b)} $f=y\,e^{5x+2y^{2}}$.

For $f_x$, the factor $y$ is a constant and the chain rule gives
$\tfrac{\partial}{\partial x}(5x+2y^{2})=5$:
\[
f_x=5y\,e^{5x+2y^{2}} .
\]
For $f_y$, both the leading $y$ and the exponent depend on $y$, so the product
rule is needed, with $\tfrac{\partial}{\partial y}(5x+2y^{2})=4y$:
\[
f_y=(1)e^{5x+2y^{2}}+y(4y)e^{5x+2y^{2}}
=\left(1+4y^{2}\right)e^{5x+2y^{2}} .
\]
\[
f_x(3,4)=5(4)e^{15+32}=20e^{47},\qquad
f_y(1,2)=(1+16)e^{5+8}=17e^{13}\approx 7.55\times10^{6}.
\]

\medskip
\textbf{(c)} $f=(x+1)^{2}+(y-3)^{3}+5xy^{3}-2$. Term by term:
\[
f_x=2(x+1)+5y^{3},\qquad
f_y=3(y-3)^{2}+15xy^{2} .
\]
The term $(y-3)^{3}$ vanishes under $\partial/\partial x$ because it contains no
$x$; likewise $(x+1)^{2}$ vanishes under $\partial/\partial y$.
\[
f_x(3,4)=2(4)+5(64)=8+320=328,
\]
\[
f_y(1,2)=3(2-3)^{2}+15(1)(4)=3+60=63 .
\]

\ans{(a) $f_x=\dfrac{(1+3y^{2})(y^{2}-2x^{3})}{(x^{3}+y^{2})^{2}}$,
$f_y=\dfrac{2xy(3x^{3}-1)}{(x^{3}+y^{2})^{2}}$;
$f_x(3,4)=-\tfrac{1862}{1849}\approx-1.007$, $f_y(1,2)=\tfrac{8}{25}=0.32$.

\smallskip
(b) $f_x=5y\,e^{5x+2y^{2}}$, $f_y=(1+4y^{2})e^{5x+2y^{2}}$;
$f_x(3,4)=20e^{47}$, $f_y(1,2)=17e^{13}$.

\smallskip
(c) $f_x=2(x+1)+5y^{3}$, $f_y=3(y-3)^{2}+15xy^{2}$;
$f_x(3,4)=328$, $f_y(1,2)=63$.}
\end{solution}

\begin{pitfall}
In (b), $f_y$ needs the product rule because $y$ appears twice --- once as a
coefficient and once inside the exponent. Writing $f_y=4y^{2}e^{5x+2y^{2}}$
drops the term from differentiating the leading $y$.
\end{pitfall}

%---------------------------------------------------------------------
\section{Verifying Laplace's equation}

\begin{question}{Q.2}
If $z=x^{2}-y^{2}+2xy$, show that
$\dfrac{\partial^{2}z}{\partial x^{2}}+\dfrac{\partial^{2}z}{\partial y^{2}}=0$.
\end{question}

\begin{solution}
Differentiate twice with respect to each variable in turn.

\medskip
With respect to $x$ (treating $y$ as constant, so $-y^{2}$ vanishes and $2xy$
gives $2y$):
\[
\pd{z}{x}=2x+2y
\;\Longrightarrow\;
\pdd{z}{x}=2 .
\]

With respect to $y$ (treating $x$ as constant, so $x^{2}$ vanishes and $2xy$
gives $2x$):
\[
\pd{z}{y}=-2y+2x
\;\Longrightarrow\;
\pdd{z}{y}=-2 .
\]

Adding:
\[
\pdd{z}{x}+\pdd{z}{y}=2+(-2)=0 . \qquad\blacksquare
\]

\medskip
A function satisfying this is called \emph{harmonic}, and the equation is
Laplace's equation. Note that the cross term $2xy$ contributes nothing to either
second derivative --- it is linear in each variable separately, so it
differentiates away.

\ans{$z_{xx}=2$ and $z_{yy}=-2$, so $z_{xx}+z_{yy}=0$ as required.}
\end{solution}

%---------------------------------------------------------------------
\section{Higher-order mixed partials}

\begin{question}{Q.3}
Let $f(x,y,z)=x^{3}y^{5}z^{7}+xy^{2}+y^{3}z-x^{2}z^{2}$. Find
(a) $f_{xz}$, (b) $f_{yz}$, (c) $f_{xy}$, (d) $f_{xyz}$, (e) $f_{yzx}$,
(f) $f_{xzy}$, (g) $f_{xxy}$, (h) $f_{xxz}$, (i) $f_{yyx}$, (j) $f_{yzy}$.
\end{question}

\begin{solution}
Compute the three first partials once, then reuse them throughout.
\[
\begin{aligned}
f_x&=3x^{2}y^{5}z^{7}+y^{2}-2xz^{2},\\
f_y&=5x^{3}y^{4}z^{7}+2xy+3y^{2}z,\\
f_z&=7x^{3}y^{5}z^{6}+y^{3}-2x^{2}z .
\end{aligned}
\]

\begin{center}
\small\sffamily
\begin{tabular}{@{}c L{5.0cm} L{6.6cm}@{}}
\toprule
 & Differentiate & Result \\
\midrule
(a) $f_{xz}$   & $f_x$ w.r.t.\ $z$ & $21x^{2}y^{5}z^{6}-4xz$ \\
(b) $f_{yz}$   & $f_y$ w.r.t.\ $z$ & $35x^{3}y^{4}z^{6}+3y^{2}$ \\
(c) $f_{xy}$   & $f_x$ w.r.t.\ $y$ & $15x^{2}y^{4}z^{7}+2y$ \\
(d) $f_{xyz}$  & $f_{xy}$ w.r.t.\ $z$ & $105x^{2}y^{4}z^{6}$ \\
(e) $f_{yzx}$  & $f_{yz}$ w.r.t.\ $x$ & $105x^{2}y^{4}z^{6}$ \\
(f) $f_{xzy}$  & $f_{xz}$ w.r.t.\ $y$ & $105x^{2}y^{4}z^{6}$ \\
(g) $f_{xxy}$  & $f_{xx}=6xy^{5}z^{7}-2z^{2}$, then $y$ & $30xy^{4}z^{7}$ \\
(h) $f_{xxz}$  & $f_{xx}$ w.r.t.\ $z$ & $42xy^{5}z^{6}-4z$ \\
(i) $f_{yyx}$  & $f_{yy}=20x^{3}y^{3}z^{7}+2x+6yz$, then $x$ & $60x^{2}y^{3}z^{7}+2$ \\
(j) $f_{yzy}$  & $f_{yz}$ w.r.t.\ $y$ & $140x^{3}y^{3}z^{6}+6y$ \\
\bottomrule
\end{tabular}
\end{center}

\vcheck{Parts (d), (e) and (f) are the same third-order derivative taken in
three different orders, and all three give $105x^{2}y^{4}z^{6}$ --- exactly as
Clairaut's theorem predicts. If any one of them had differed, an arithmetic
error would have been located.}

\ans{As tabulated. In particular
$f_{xyz}=f_{yzx}=f_{xzy}=105x^{2}y^{4}z^{6}$.}
\end{solution}

%---------------------------------------------------------------------
\section{Interpreting partial derivatives of a cost function}

\begin{question}{Q.4}
Suppose $z=f(x,y)$ represents the cost of constructing a certain structure,
where $x$ and $y$ express the labour and material cost respectively. Explain
what $f_x$ and $f_y$ represent.
\end{question}

\begin{solution}
Both are \emph{marginal costs} --- rates of change of total cost --- but each
with respect to one input while the other is frozen.

\begin{description}
  \item[$f_x$: marginal cost with respect to labour.] The rate at which total
        construction cost changes per one-unit increase in labour cost, with
        material cost held fixed. If $f_x=1.4$, then spending one more rupee on
        labour raises the total cost of the structure by about Rs.~1.40 ---
        more than the rupee itself, indicating that labour drags other costs up
        with it.
  \item[$f_y$: marginal cost with respect to material.] The rate at which total
        cost changes per one-unit increase in material cost, with labour held
        fixed.
\end{description}

The word ``held fixed'' carries the meaning. In practice labour and material
costs move together, so $f_x$ answers a deliberately artificial but useful
question: \emph{if only labour changed, what would happen to the total?} That
isolation is what makes the partial derivative the right tool for attributing a
cost overrun to one input rather than another.

\ans{$f_x$ is the marginal cost of the structure with respect to labour --- the
change in total cost per unit change in labour cost, material held constant.
$f_y$ is the corresponding marginal cost with respect to material, labour held
constant.}
\end{solution}

%---------------------------------------------------------------------
\section{Relative extrema of functions of two variables}

\begin{question}{Q.6}
Find all points where the following functions have relative extrema.
\textbf{(a)} $f(x,y)=(y^{2}-4)(e^{x}-1)$ \quad
\textbf{(b)} $f(x,y)=x^{2}+y^{2}+xy-9x+1$ \quad
\textbf{(c)} $f(x,y)=4xy-x^{4}-y^{4}$ \quad
\textbf{(d)} $f(x,y)=x^{2}+y^{2}+\dfrac{32}{xy}$ \quad
\textbf{(e)} $f(x,y)=x\,e^{y+1}$
\end{question}

\begin{solution}
Method throughout: set $f_x=f_y=0$ to locate critical points, then classify each
with $D=f_{xx}f_{yy}-(f_{xy})^{2}$.

\medskip
\textbf{(a)} $f_x=(y^{2}-4)e^{x}$ and $f_y=2y(e^{x}-1)$.

Since $e^{x}>0$ always, $f_x=0$ forces $y^{2}=4$, i.e.\ $y=\pm2$. Substituting
into $f_y=0$: $2(\pm2)(e^{x}-1)=0$ requires $e^{x}=1$, so $x=0$. Critical
points: $(0,2)$ and $(0,-2)$.
\[
f_{xx}=(y^{2}-4)e^{x},\qquad f_{yy}=2(e^{x}-1),\qquad f_{xy}=2ye^{x}.
\]
At both points $f_{xx}=0$ and $f_{yy}=0$, while $f_{xy}=\pm4$:
\[
D=0\cdot 0-(\pm4)^{2}=-16<0 .
\]
\ans{(a) Critical points $(0,2)$ and $(0,-2)$; both are \textbf{saddle points}
($D=-16<0$). The function has no relative extrema.}

\medskip
\textbf{(b)} $f_x=2x+y-9$ and $f_y=2y+x$.

From $f_y=0$, $x=-2y$. Substituting into $f_x=0$:
\[
2(-2y)+y-9=0 \;\Longrightarrow\; -3y=9 \;\Longrightarrow\; y=-3,\quad x=6 .
\]
\[
f_{xx}=2,\quad f_{yy}=2,\quad f_{xy}=1,
\qquad D=(2)(2)-1^{2}=3>0,\quad f_{xx}=2>0 .
\]
\[
f(6,-3)=36+9+6(-3)-54+1=-26 .
\]
\ans{(b) Relative \textbf{minimum} at $(6,-3)$, value $f=-26$.}

\medskip
\textbf{(c)} $f_x=4y-4x^{3}$ and $f_y=4x-4y^{3}$.

Setting both to zero gives $y=x^{3}$ and $x=y^{3}$. Substituting the first into
the second:
\[
x=(x^{3})^{3}=x^{9} \;\Longrightarrow\; x^{9}-x=0
\;\Longrightarrow\; x(x^{8}-1)=0 ,
\]
so $x=0,\,1,\,-1$ (the real roots), giving critical points $(0,0)$, $(1,1)$,
$(-1,-1)$.
\[
f_{xx}=-12x^{2},\quad f_{yy}=-12y^{2},\quad f_{xy}=4,
\qquad D=144x^{2}y^{2}-16 .
\]

\begin{center}
\small\sffamily
\begin{tabular}{@{}C{2.4cm} C{2.2cm} C{2.2cm} C{3.2cm} C{2.0cm}@{}}
\toprule
Point & $D$ & $f_{xx}$ & Classification & $f$ \\
\midrule
$(0,0)$   & $-16$ & $0$   & saddle point      & $0$ \\
$(1,1)$   & $128$ & $-12$ & relative maximum  & $2$ \\
$(-1,-1)$ & $128$ & $-12$ & relative maximum  & $2$ \\
\bottomrule
\end{tabular}
\end{center}
\ans{(c) Relative \textbf{maxima} at $(1,1)$ and $(-1,-1)$, both of value 2;
$(0,0)$ is a saddle point.}

\medskip
\textbf{(d)} Write $f=x^{2}+y^{2}+32x^{-1}y^{-1}$, so
\[
f_x=2x-\frac{32}{x^{2}y},\qquad f_y=2y-\frac{32}{xy^{2}} .
\]
Setting each to zero:
\[
2x^{3}y=32 \;\Rightarrow\; x^{3}y=16,
\qquad
2xy^{3}=32 \;\Rightarrow\; xy^{3}=16 .
\]
Dividing one by the other gives $x^{2}/y^{2}=1$, so $y=\pm x$.

If $y=x$: $x^{4}=16 \Rightarrow x=\pm2$, giving $(2,2)$ and $(-2,-2)$.
If $y=-x$: $-x^{4}=16$, which has no real solution.
\[
f_{xx}=2+\frac{64}{x^{3}y},\qquad f_{yy}=2+\frac{64}{xy^{3}},\qquad
f_{xy}=\frac{32}{x^{2}y^{2}} .
\]
At $(2,2)$: $f_{xx}=2+\tfrac{64}{16}=6$, $f_{yy}=6$, $f_{xy}=\tfrac{32}{16}=2$,
\[
D=36-4=32>0,\quad f_{xx}=6>0 \;\Longrightarrow\; \text{minimum},
\]
\[
f(2,2)=4+4+\frac{32}{4}=16 .
\]
At $(-2,-2)$ the same values arise ($x^{3}y=16$ and $x^{2}y^{2}=16$ again), so
$D=32>0$, $f_{xx}=6>0$, and $f(-2,-2)=16$.
\ans{(d) Relative \textbf{minima} at $(2,2)$ and $(-2,-2)$, both of value 16.}

\medskip
\textbf{(e)} $f_x=e^{y+1}$ and $f_y=xe^{y+1}$.

Since $e^{y+1}>0$ for every $y$, $f_x$ is \emph{never} zero. There are no
critical points, hence nothing to classify.
\ans{(e) No critical points ($f_x=e^{y+1}\neq0$ everywhere), so the function has
\textbf{no relative extrema}.}
\end{solution}

\begin{pitfall}
$D<0$ means saddle point --- a genuine, complete answer, not a failure. And
$D=0$ means the test is \emph{inconclusive}, not that the point is a saddle;
in that case the test simply cannot decide and other reasoning is needed.
\end{pitfall}

%---------------------------------------------------------------------
\section{A production function --- and a second defective question}
\label{sec:u8q7}

\begin{question}{Q.7}
Suppose a production function $P$ is given by
\[
P=f(l,k)=0.8l^{3}-0.02l^{2}+2.5k^{2}-0.05k^{3}.
\]
Find the values of $l$ and $k$ so as to maximise output $P$.
\end{question}

\begin{solution}
\textbf{As printed, no maximum exists.} Take the partials:
\[
\pd{P}{l}=2.4l^{2}-0.04l=l(2.4l-0.04),\qquad
\pd{P}{k}=5k-0.15k^{2}=k(5-0.15k).
\]
Setting them to zero gives $l=0$ or $l=\tfrac{0.04}{2.4}=\tfrac{1}{60}$, and
$k=0$ or $k=\tfrac{5}{0.15}=\tfrac{100}{3}$.

Now the second-order condition in $l$:
\[
\pdd{P}{l}=4.8l-0.04, \qquad
\text{at } l=\tfrac{1}{60}: \quad 4.8\left(\tfrac{1}{60}\right)-0.04
=0.08-0.04=0.04>0 .
\]
The critical point in $l$ is a \emph{minimum}, not a maximum. And more
decisively, the leading term $0.8l^{3}$ is positive and cubic, so
\[
P\to +\infty \quad\text{as}\quad l\to\infty .
\]
An unbounded function has no maximum. The question as printed cannot be
answered.

\ans{As printed, $P$ is unbounded above in $l$ (the $+0.8l^{3}$ term dominates),
so \textbf{no maximising $(l,k)$ exists}. The stationary point
$l=\tfrac{1}{60}$ is a minimum in $l$, since $P_{ll}=0.04>0$ there.}

\medskip
\textbf{The intended problem.} A production function that can be maximised needs
a \emph{negative} cubic term, and the standard form of this exercise is
\[
P=f(l,k)=0.8l^{2}-0.02l^{3}+2.5k^{2}-0.05k^{3},
\]
where the exponents on the first two terms are interchanged. Then
\[
\pd{P}{l}=1.6l-0.06l^{2}=l(1.6-0.06l)=0
\;\Longrightarrow\; l=0 \ \text{ or } \ l=\frac{1.6}{0.06}=\frac{80}{3}\approx 26.67,
\]
\[
\pd{P}{k}=5k-0.15k^{2}=k(5-0.15k)=0
\;\Longrightarrow\; k=0 \ \text{ or } \ k=\frac{5}{0.15}=\frac{100}{3}\approx 33.33 .
\]
Classifying the non-trivial point:
\[
\pdd{P}{l}=1.6-0.12l \;\Big|_{\,l=80/3}=1.6-3.2=-1.6<0,
\]
\[
\pdd{P}{k}=5-0.3k \;\Big|_{\,k=100/3}=5-10=-5<0,
\qquad
P_{lk}=0 ,
\]
\[
D=(-1.6)(-5)-0^{2}=8>0,\quad P_{ll}<0 \;\Longrightarrow\;
\text{relative maximum}.
\]
Maximum output:
\[
P=0.8\left(\tfrac{80}{3}\right)^{2}-0.02\left(\tfrac{80}{3}\right)^{3}
+2.5\left(\tfrac{100}{3}\right)^{2}-0.05\left(\tfrac{100}{3}\right)^{3}
\approx 189.63+925.93=1115.56 .
\]
\ans{For the intended function $P=0.8l^{2}-0.02l^{3}+2.5k^{2}-0.05k^{3}$:
output is maximised at $l=\tfrac{80}{3}\approx 26.67$ and
$k=\tfrac{100}{3}\approx 33.33$, with $D=8>0$ and $P_{ll}=-1.6<0$ confirming a
maximum; $P_{\max}\approx 1115.6$.}
\end{solution}

\begin{pitfall}
Never stop at ``$P_l=0$ and $P_k=0$, therefore maximum''. The second-order test
is what distinguishes a maximum from a minimum or a saddle, and in this question
it is what reveals that the printed function has no maximum at all. Examiners
award the second-order working separately.
\end{pitfall}

%---------------------------------------------------------------------
\section{Maximising profit in two variables}

\begin{question}{Q.8}
The profit of a certain company is given by
\[
P(x,y)=1200+80x-2x^{2}+100y-y^{2},
\]
where $x$ and $y$ represent the cost of a unit of labour and a unit of goods
respectively. Find $x$ and $y$ that maximise the profit, and find the maximum
profit.
\end{question}

\begin{solution}
The variables separate --- there is no $xy$ term --- so each can be optimised
independently. Setting the first partials to zero:
\[
\pd{P}{x}=80-4x=0 \;\Longrightarrow\; x=20,
\]
\[
\pd{P}{y}=100-2y=0 \;\Longrightarrow\; y=50 .
\]

Second-order test:
\[
P_{xx}=-4,\qquad P_{yy}=-2,\qquad P_{xy}=0,
\]
\[
D=(-4)(-2)-0^{2}=8>0 \quad\text{and}\quad P_{xx}=-4<0
\;\Longrightarrow\; \text{relative maximum}.
\]

Maximum profit:
\[
P(20,50)=1200+80(20)-2(400)+100(50)-(50)^{2}
\]
\[
=1200+1600-800+5000-2500=4500 .
\]

\vcheck{Perturbing either variable lowers $P$: $P(21,50)=4498$ and
$P(20,51)=4499$. Consistent with a maximum at $(20,50)$.}

\ans{Profit is maximised at $x=20$ and $y=50$, giving a maximum profit of
\textbf{4{,}500}.}
\end{solution}

%---------------------------------------------------------------------
\section{Maximising revenue with an interaction term}

\begin{question}{Q.9}
The revenue in thousands of dollars from selling $x$ spas and $y$ solar heaters
is
\[
R(x,y)=15+80x+81y-2x^{2}-3y^{2}-4xy .
\]
Find the quantity of each that must be sold to generate maximum revenue, and
find the maximum revenue.
\end{question}

\begin{solution}
Here the $-4xy$ term couples the two products, so the equations must be solved
\emph{simultaneously}.
\[
\pd{R}{x}=80-4x-4y=0 \;\Longrightarrow\; x+y=20, \tag{1}
\]
\[
\pd{R}{y}=81-6y-4x=0 \;\Longrightarrow\; 4x+6y=81. \tag{2}
\]
From (1), $x=20-y$. Substituting into (2):
\[
4(20-y)+6y=81 \;\Longrightarrow\; 80+2y=81 \;\Longrightarrow\; y=0.5,
\]
\[
x=20-0.5=19.5 .
\]

Second-order test:
\[
R_{xx}=-4,\qquad R_{yy}=-6,\qquad R_{xy}=-4,
\]
\[
D=(-4)(-6)-(-4)^{2}=24-16=8>0,\quad R_{xx}=-4<0
\;\Longrightarrow\; \text{maximum}.
\]

Maximum revenue:
\[
\begin{aligned}
R(19.5,\,0.5)&=15+80(19.5)+81(0.5)-2(19.5)^{2}-3(0.5)^{2}-4(19.5)(0.5)\\
&=15+1560+40.5-760.5-0.75-39\\
&=815.25 .
\end{aligned}
\]
Revenue is measured in thousands of dollars, so this is \$815{,}250.

\medskip
\textbf{A practical note.} Spas and heaters are sold in whole units, so the
exact optimum $(19.5,\,0.5)$ is not attainable. Testing the four nearest
integer points:
\[
R(20,0)=815.00,\quad R(19,1)=815.00,\quad R(20,1)=813.00,\quad R(19,0)=813.00 .
\]
The best integer plans are 20 spas with no heaters, or 19 spas with 1 heater,
each giving \$815{,}000 --- \$250 below the continuous optimum.

\ans{Revenue is maximised at $x=19.5$ spas and $y=0.5$ heaters, with
$R_{\max}=815.25$ thousand dollars ($\$815{,}250$). Restricted to whole units,
$(20,0)$ or $(19,1)$ both give \$815{,}000.}
\end{solution}

\begin{pitfall}
With a cross term present, $D=f_{xx}f_{yy}-(f_{xy})^{2}$ can go negative even
when both $f_{xx}$ and $f_{yy}$ are negative. Here $24-16=8$ only just clears
zero. Never skip the $(f_{xy})^{2}$ subtraction on the grounds that both second
derivatives ``look right''.
\end{pitfall}
